\chapter{Conclusion}

\section{Contributions}
This work characterised the migration of the data contracts of a real MCP server to an
ontology. It delivered an ontology whose content is generated from the 80 contracts, a
ledger that states how each contract construct was represented, an evaluation with a
deterministic router and with an LLM-based agent on questions written by people, and a
prototype server that answers from the ontology.

The main findings are three. More than half of the contract content becomes typed knowledge
and 40\% survives as literals; what is left behind is physical and operational. For the
lexical router of the server the ontology is at best equivalent to the contracts, and loses
5 to 8 points of top-1 accuracy when free text is removed. For an LLM-based agent the
ontology with typed content only is as accurate as the contracts and uses between a third
and almost three fifths less context, under three models of different sizes.

\section{Limitations}
The study covers one organisation, one catalogue and one language. Effectiveness was
measured as the choice of the dataset, not as the correctness of the final answer. The agent
round used three models of one family, and each condition was run once on non-deterministic
models; a repeat of two conditions changed the outcome of three to four questions out of 106. The ontology content derives from
the contracts, so the study says little about ontologies built from independent domain
knowledge. The question sets are not public.

\section{Future Work}
\begin{itemize}[leftmargin=1.5cm]
\item Evaluate the final answer: execute the generated SQL and compare values, including
executable but wrong answers caused by units and aggregation.
\item Add knowledge that the contracts do not hold, such as the relations between plants,
reservoirs and basins, and measure questions that need it.
\item Obtain a second annotation of the question sets and report inter-rater agreement.
\item Replace the lexical router with a retriever that uses the ontology graph.
\item Replicate the migration on the catalogue of another organisation.
\end{itemize}

\section{Lessons Learned}
The two attempts at the migration taught more than the final numbers.

\textbf{Generate, do not curate.} The first attempt tuned hand-written aliases until the
coverage of the datasets looked good, which fitted the ontology to the test. Deriving every
instance from the contracts with a program made the mapping repeatable, auditable through
the ledger, and free of parameters fitted to the questions.

\textbf{Evaluate on questions written by people, and check them.} The first test set was
made of automatic paraphrases of the development questions and overstated the results.
The sets written by ONS staff changed the conclusions; they also overlap and disagree on
four questions, which only a comparison between them revealed.

\textbf{Correct for multiple comparisons and repeat the runs.} A difference that looked
significant in the first attempt disappeared after the Holm correction, and repeating the
agent round showed that the variation between runs is as large as the difference between
contracts and ontology.

\textbf{Measure cost together with accuracy.} The main finding of this study, the smaller
context of the ontology catalogue, would have been missed if only accuracy had been
measured.

\textbf{Verify what assistants and sources claim.} Several statements produced during the
work, by AI assistants and by the press, did not hold when checked against the code, the
data or primary sources. Version control from the first day, which the first attempt
lacked, is what made these checks possible.

If the work were started again, a domain expert would review the schema of the ontology
and the mapping decisions before the experiments, a condition with the contracts stripped
of free text would separate the effect of the ontology from that of a shorter description,
and the agent would be evaluated against the real server and on the final answer.
